% Copyright (c) 2025 ZKLib Contributors. All rights reserved.
% Released under Apache 2.0 license as described in the file LICENSE.
% Authors: Poulami Das (Least Authority)

\section{WHIR}\label{sec:whir}

WHIR is the Reed--Solomon IOPP layer whose soundness depends on three families of
formalized prerequisites: mutual correlated agreement for random linear combinations, WHIR
folding and block-distance list decoding, and out-of-domain sampling.  ArkLib now exposes the
paper theorem as named front doors and separates the remaining mathematical residuals from the
protocol scaffold.  The genuine open kernel is the Johnson-regime MCA/capture argument; the
protocol objects, completeness shells, budget accounting, and conditional RBR entry points are
represented explicitly.  The live implementation tracker is
\href{https://github.com/lalalune/ArkLib/issues/302}{WHIR issue \#302}; this blueprint records
that reduced state without upgrading the conditional Lean front doors into unconditional theorem
closure.

\subsection{Reed--Solomon proximity and MCA}\label{sec:whir_mca}

\begin{definition}[Proximity generator]\label{def:whir_proximity_generator}
\lean{Generator.ProximityGenerator}
\uses{def:linear_code,def:distance_from_code}
    A generator $\mathsf{Gen}$ for a linear code $\code$ is a proximity generator with
    bound $B$ and error $\epsilon$ when the following implication holds: if a random linear
    combination of words is close to $\code$ with probability above $\epsilon$, then all input
    words agree with codewords of $\code$ on one common large coordinate set.
\end{definition}

\begin{theorem}[Reed--Solomon proximity generator]\label{thm:whir_proximity_gap}
\lean{RSGenerator.genRSC}
\uses{def:whir_proximity_generator,def:reed_solomon_code}
    For a Reed--Solomon code $\code=\rscode[\field,\evaldomain,m]$, the generator
    $(1,\alpha,\ldots,\alpha^{\parl-1})$ is the WHIR proximity generator with Johnson-radius
    proximity bound and the corresponding finite-field error term.
\end{theorem}

\begin{definition}[Mutual correlated agreement]\label{def:whir_mca}
\lean{MutualCorrAgreement.hasMutualCorrAgreement}
\uses{def:whir_proximity_generator}
    Mutual correlated agreement strengthens the proximity-generator conclusion: except with
    the stated error, a random linear combination cannot be close to $\code$ on a large set
    while some input word fails to agree with any codeword on that same set.
\end{definition}

\begin{lemma}[MCA from proximity generation]\label{lem:whir_mca_from_proximity}
\lean{MutualCorrAgreement.mca_linearCode,MutualCorrAgreement.mca_linearCode_udrFree}
\uses{def:whir_proximity_generator,def:whir_mca}
    A proximity generator for a linear code yields a mutual-correlated-agreement generator,
    with the proximity radius capped by the unique-decoding radius of the code.  The UDR-free
    variant packages the same implication in the form used by later WHIR bridges.
\end{lemma}

\begin{theorem}[WHIR MCA front doors]\label{thm:whir_mca_front_doors}
\lean{MutualCorrAgreement.mca_rsc,MutualCorrAgreement.mca_johnson_bound_CONJECTURE,MutualCorrAgreement.mca_capacity_bound_CONJECTURE,MutualCorrAgreement.mca_johnson_bound_CONJECTURE_pair_of_johnsonNumericBound,MutualCorrAgreement.mca_johnson_bound_CONJECTURE_pair_of_claim1_cells,MutualCorrAgreement.mca_johnson_bound_CONJECTURE_pair_via_ellary,MutualCorrAgreement.mca_johnson_bound_CONJECTURE_smallField}
\uses{def:whir_mca,def:smooth_rs_code,thm:whir_proximity_gap}
    ArkLib records the WHIR MCA target for smooth Reed--Solomon codes up to the Johnson bound
    and retains the capacity-shaped statement as a historical front door.  The current
    formalization also exposes pair, curve, Ellary, and small-field routes that reduce the
    Johnson theorem to explicit numeric or capture hypotheses instead of hiding those
    assumptions inside the protocol theorem.
\end{theorem}

\begin{remark}[WHIR MCA status]\label{rem:whir_mca_status}
\lean{MutualCorrAgreement.hasMutualCorrAgreement_genRSC_pair_of_epsMCA_le,MutualCorrAgreement.hasMutualCorrAgreement_genRSC_of_epsMCACurve_le,Whir302B.hasMutualCorrAgreement_genRSC_of_epsMCAP_le,Whir302B.hasMutualCorrAgreement_genRSC_vandermonde_of_epsMCAP_le,MutualCorrAgreement.hasMutualCorrAgreement_genRSC_pair_of_johnsonNumericBound,CodingTheory.ProximityGap.Hab25Core.Hab25JohnsonEndgame.johnsonNumericBound_of_affine_capture}
\uses{thm:whir_mca_front_doors}
    The Johnson-radius statement is the live mathematical target.  The affine-line pair
    generator seam, higher-arity curve/parameter-sampling seams, the Johnson-numeric bridge,
    and small-field/vacuous exits are formalized.  What remains is the affine-capture or
    generic-to-scalar decoding kernel: bad scalar-close Reed--Solomon codewords must be shown
    to come from the bounded list of generic $\field(Z)$-linear factors rather than from
    higher-degree factors after specialization.
\end{remark}

\begin{lemma}[MCA bridge for folded list decoding]\label{lem:whir_mca_list_decoding}
\lean{Fold.folding_preserves_listdecoding_base_of_mca_bridge}
\uses{def:whir_mca,def:interleaved_code,def:list_close_codewords}
    The current production bridge threads a genuine \texttt{hasMutualCorrAgreement}
    hypothesis into the one-step folded list-decoding bound.  Together with the deterministic
    forward-inclusion half, it gives the reverse-inclusion probability estimate consumed by
    WHIR's folding-preserves-list-decoding theorem.
\end{lemma}

\subsection{Folding and block distance}\label{sec:whir_folding}

\begin{definition}[Square-root extraction]\label{def:whir_extract}
\lean{Fold.extract_x}
    The WHIR fold queries a word on paired points.  The extractor chooses a square root
    representative of a point in the next folded domain so that both preimages can be queried.
\end{definition}

\begin{definition}[WHIR folding]\label{def:whir_folding}
\lean{Fold.foldf,Fold.fold_k_core,Fold.fold_k,Fold.fold_k_set}
\uses{def:whir_extract}
    The one-step fold maps $f$ to
    \[
        y \mapsto {f(x)+f(-x)\over 2}
        + \alpha {f(x)-f(-x)\over 2x},
    \]
    where $x$ is extracted from $y$.  Iterating this operation gives \texttt{fold\_k};
    applying it pointwise to a set of words gives \texttt{fold\_k\_set}.
\end{definition}

\begin{lemma}[Polynomial semantics of folding]\label{lem:whir_folding_polynomial}
\lean{Fold.foldPoly,Fold.foldf_eq_foldPoly_eval,Fold.foldf_step_mem_smoothCode,Fold.fold_f_g_core,Fold.fold_f_g,Fold.fold_f_g_poly}
\uses{def:whir_folding,def:smooth_rs_code}
    The functional fold agrees with the polynomial even/odd decomposition, preserves smooth
    Reed--Solomon membership under the expected degree conditions, and specializes the folded
    polynomial to the multilinear-extension view used in WHIR.
\end{lemma}

\begin{definition}[Block relative distance]\label{def:whir_block_distance}
\lean{BlockRelDistance.block,BlockRelDistance.blockRelDistance,BlockRelDistance.minBlockRelDistance,BlockRelDistance.listBlockRelDistance}
\uses{def:smooth_rs_code}
    Block distance measures disagreement after grouping coordinates into folding blocks.  The
    associated block list contains codewords within a given block-relative radius from the
    received word.
\end{definition}

\begin{lemma}[Block distance bounds Hamming distance]\label{lem:whir_block_distance_bounds}
\lean{BlockRelDistance.blockRelDistance_eq_relHammingDist_of_k_eq_i,BlockRelDistance.relHammingDist_le_blockRelDistance,BlockRelDistance.listBlock_subset_listHamming}
\uses{def:whir_block_distance,def:list_close_codewords}
    At zero fold depth, block distance is ordinary relative Hamming distance.  In general,
    Hamming distance is bounded by block distance, so block-close lists are subsets of ordinary
    close-codeword lists.
\end{lemma}

\begin{theorem}[Folding preserves list decoding]\label{thm:whir_folding_preserves_list_decoding}
\lean{Fold.folding_listdecoding_if_genMutualCorrAgreement,Fold.folding_preserves_listdecoding_base,Fold.folding_preserves_listdecoding_base_of_mca_bridge,Fold.folding_preserves_listdecoding_bound,Fold.folding_preserves_listdecoding_base_ne_subset}
\uses{def:whir_folding,def:whir_block_distance,lem:whir_mca_list_decoding,lem:whir_folding_polynomial}
    Assuming MCA for the relevant Reed--Solomon codes, a random WHIR fold preserves the
    intended list-decoding invariant: the fold of the block-close list agrees with the
    close-codeword list of the folded word except with the accumulated MCA and degree-error
    budget.
\end{theorem}

\subsection{Out-of-domain sampling}\label{sec:whir_ood}

\begin{lemma}[Constrained and unconstrained agreement equivalence]\label{lem:whir_ood_equiv}
\lean{OutOfDomSmpl.out_of_dom_smpl_1}
\uses{def:smooth_rs_code,def:list_close_codewords}
    Agreement of two close Reed--Solomon codewords on sampled out-of-domain points is equivalent
    to non-uniqueness in the corresponding constrained smooth Reed--Solomon list.
\end{lemma}

\begin{lemma}[Out-of-domain sampling bound]\label{lem:whir_ood_bound}
\lean{OutOfDomSmpl.out_of_dom_smpl_2}
\uses{lem:whir_ood_equiv,def:list_decodable}
    If the base Reed--Solomon code is list decodable, repeated out-of-domain sampling bounds
    the probability that a constrained list has more than one candidate by the usual
    list-size-squared collision term.
\end{lemma}

\subsection{Round-by-round WHIR soundness}\label{sec:whir_rbr}

\begin{definition}[WHIR parameters and relation]\label{def:whir_params}
\lean{WhirIOP.Params,WhirIOP.ParamConditions,Fold.GenMutualCorrParams,WhirIOP.OracleStatement,WhirIOP.whirRelation}
\uses{def:smooth_rs_code,def:list_decodable}
    The WHIR parameter structures store the round count, folding schedule, evaluation domains,
    repetition parameters, distance targets, and MCA/list-decoding assumptions.  The relation
    \texttt{whirRelation} is the initial constrained smooth Reed--Solomon proximity statement.
\end{definition}

\begin{theorem}[WHIR round-by-round soundness front door]\label{thm:whir_rbr_soundness}
\lean{WhirIOP.whir_rbr_soundness}
\uses{def:whir_params,thm:whir_folding_preserves_list_decoding,lem:whir_ood_bound,thm:whir_mca_front_doors}
    The paper-level RBR theorem asserts existence of a vector IOPP for the WHIR relation with
    the stated fold, out-of-domain, shift, and final-round error budgets.  In Lean this is the
    front-door theorem consumed by protocol-level constructions once the real security gap and
    MCA hypotheses are supplied.
\end{theorem}

\begin{definition}[WHIR transcript vector IOP]\label{def:whir_transcript_iop}
\lean{WhirIOP.Construction.whirPaperTranscriptVectorSpec,whirMakeTranscript,whirVerify,whirVectorIOP,whirVectorIOP_perfectCompleteness,whirVectorIOP_rbrKnowledgeSoundness,whirVectorIOP_isSecureWithGap}
\uses{def:whir_params}
    The protocol layer builds the paper transcript shape, prover transcript construction,
    verifier, and vector IOP wrapper.  The current verifier shell has a pure-true placeholder
    acceptance predicate, so its unconditional soundness surface is intentionally separated from
    the intended small-error theorem.
\end{definition}

\begin{theorem}[WHIR protocol completeness]\label{thm:whir_protocol_completeness}
\lean{Whir302.whirVectorIOP_perfectCompleteness_holds,Whir302.whirVectorIOP_isSecureWithGap_of_rbr,whirVectorIOP_isSecureWithGap_holds}
\uses{def:whir_transcript_iop}
    The assembled WHIR vector IOP is perfectly complete for the current protocol scaffold, and
    a supplied RBR knowledge-soundness proof instantiates the secure-with-gap front door.
\end{theorem}

\begin{theorem}[WHIR dummy soundness surface]\label{thm:whir_dummy_soundness}
\lean{WhirIOP.whirVectorIOP_rbrKnowledgeSoundness_dummy,WhirIOP.whirVectorIOP_isSecureWithGap_dummy_of_rbr}
\uses{def:whir_transcript_iop}
    The dummy verifier proves only the trivial $\varepsilon_{\mathrm{rbr}}=1$ soundness
    surface.  This theorem is useful as plumbing, but it is not the WHIR paper soundness claim.
\end{theorem}

\begin{definition}[RBR budget accounting]\label{def:whir_rbr_budget}
\lean{Issue113WHIR.maxFolds,Issue113WHIR.rbrBudgetSet,Issue113WHIR.epsRbr}
\uses{def:whir_params}
    The RBR accounting layer takes the per-round fold, out-of-domain, shift, and final errors
    and packages their least upper bound as the WHIR paper budget.
\end{definition}

\begin{theorem}[RBR budget keystones]\label{thm:whir_rbr_budget_accounting}
\lean{Issue113WHIR.rbrBudgetSet_nonempty,Issue113WHIR.epsRbr_le_of_forall_le,Issue113WHIR.epsRbr_isLUB,Issue113WHIR.epsRbr_dominates,Issue113WHIR.soundOk_epsRbr_of_keystone,Issue113WHIR.soundOk_of_keystone_of_forall_le}
\uses{def:whir_rbr_budget}
    The named accounting theorems prove nonemptiness of the budget set, least-upper-bound
    properties for $\epsilon_{\mathrm{rbr}}$, dominance over the component errors, and the
    monotone transport from the named keystone to any looser per-challenge error budget.
\end{theorem}

\begin{theorem}[Vector IOP breakthrough front doors]\label{thm:whir_vector_iop_breakthrough}
\lean{WhirIOP.Construction.whirVectorSpec,WhirIOP.Construction.whirVectorSpec_card_challengeIdx,WhirIOP.Construction.whir_rbr_soundness_of_secure_gap,WhirIOP.Construction.whir_rbr_soundness_of_whirVectorSpec_secure_gap,WhirIOPP.whir_vector_iop_breakthrough}
\uses{thm:whir_rbr_soundness,def:whir_transcript_iop,thm:whir_rbr_budget_accounting}
    Given a real secure-with-gap proof for the WHIR vector spec and the paper budget
    inequalities, the construction layer instantiates the paper RBR theorem and the
    high-level vector IOP breakthrough statement.
\end{theorem}

\begin{remark}[Current WHIR status]\label{rem:whir_status}
\lean{Issue113WHIR.WhirRbrKeystone,Issue113WHIR.whirRbrShape,Issue113WHIR.whirRbrShape_of_secure,WhirIOP.Construction.whir_rbr_soundness_of_secure_gap,WhirIOPP.whir_vector_iop_breakthrough}
\uses{rem:whir_mca_status,thm:whir_vector_iop_breakthrough}
    The blueprint is complete as a map of the current formalization.  WHIR has explicit Lean
    surfaces for proximity/MCA prerequisites, folding/list-decoding preservation,
    out-of-domain sampling, protocol assembly, completeness, and budget accounting.  The
    unconditional paper soundness theorem is not closed: it still depends on the real WHIR
    secure-with-gap proof, which in turn depends on the Johnson MCA capture kernel described
    above.  The dummy all-one error theorem and the conditional vector-IOP breakthrough theorem
    are audit scaffolding, not a proof of the intended sub-one RBR bound.
\end{remark}
